\subsection{Target Choice Selects Support from a Shared Projection Profile}

We next vary the scalar target while holding the hidden state fixed. This
isolates the central SRP distinction: the projection profile says which readout
directions are present in the context, while the target coefficients determine
which of those directions count as support or opposition for the selected score.

\paragraph{Target choice matters.}
A hidden state has a target-independent readout-feature projection vector; the
scalar target turns that vector into signed feature contributions. The same
hidden state can be evaluated with a row score, row-versus-reference contrast,
or pair/group contrast, so a large projection \(h^\top d_i\) has an
interpretive sign only through the readout row or contrast being scored.

Aggregate metrics match this interpretation. In the five-model Qwen/Gemma suite,
reference contrasts are easiest: 301 top-token-versus-vocabulary-mean cases
give 1,505 model--case rows with median \(\rho=0.021\) and perfect sign
agreement. Pairwise and group contrasts are harder: top-1 versus top-5 margins
have median \(\rho=0.253\), while top-1 versus top-2 margins fall to \(0.824\)
sign agreement. Accordingly, each account is interpreted at its chosen score,
with local sign and reconstruction error reported; full counts and examples are
in \Cref{tab:app-general-readout-query-families} and
Appendix~\ref{app:query-family-display-example}.
